\section*{الفصل \ref{ch:b1:spherical-astronomy} --- الفلك الكروي}

\begin{solution}{exo:b1:spherical-astronomy:1}
\omterm{def:b1:spherical-astronomy:horizontal}{ارتفاع القطب} يساوي \omterm{def:b1:spherical-astronomy:geographic}{العرض الجغرافي}: \ang{52}. لا يغيب النجم إن كان
$\delta > \ang{90} - \ang{52} = \ang{38}$، ولا يطلع إن كان
$\delta < \ang{-38}$.
\end{solution}

\begin{solution}{exo:b1:spherical-astronomy:2}
$\qty{6.75}{h}\times\ang{15} = \ang{101.25}$. $H = \mathrm{LST} - \alpha =
\qty{10}{h} - \qty{6}{h}\,\qty{45}{min} = \qty{3}{h}\,\qty{15}{min}$، أي
\ang{48.75}. \omterm{def:b1:spherical-astronomy:equatorial}{والزاوية الساعية} موجبة: فالنجم غرب خط الزوال، وكان قد عبره قبل ثلاث ساعات وربع نجمية.
\end{solution}

\begin{solution}{exo:b1:spherical-astronomy:3}
$h = \ang{90} - \varphi + \delta_\odot$: \ang{73.44} في انقلاب يونيو،
\ang{50} في الاعتدالين، \ang{26.56} في انقلاب ديسمبر.
\end{solution}

\begin{solution}{exo:b1:spherical-astronomy:4}
1900: قابلة للقسمة على 100 وغير قابلة للقسمة على 400، فليست كبيسة. 2000: قابلة للقسمة على 400، فهي كبيسة. 2100: ليست كبيسة. وتزيد السنة اليولية على \omterm{def:b1:spherical-astronomy:tropical-year}{السنة المدارية} بمقدار $365.25 - 365.2422 = 0.0078$ يوم، أي \num{0.78} يوم في القرن؛ وتزيد عليها السنة الغريغورية بمقدار $0.0003$ يوم، أي \num{0.03} يوم في القرن.
\end{solution}

\begin{solution}{exo:b1:spherical-astronomy:5}
$\cos\Delta = \sin^2\ang{45} + \cos^2\ang{45}\cos\ang{90} = 0.5$، إذن
$\Delta = \ang{60}$، والمسافة هي $6371\times\pi/3 = \qty{6672}{km}$.
على الدائرة الموازية، وهي \omterm{def:b1:spherical-astronomy:spherical-triangle}{دائرة صغرى} نصف قطرها $R\cos\ang{45}$، تكون المسافة
$6371\times\cos\ang{45}\times\pi/2 = \qty{7076}{km}$. \omterm{def:b1:spherical-astronomy:spherical-triangle}{والدائرة العظمى} أقصر بنحو
\qty{404}{km}.
\end{solution}

\begin{solution}{exo:b1:spherical-astronomy:6}
$\sin h = \sin\ang{48}\sin\ang{20} + \cos\ang{48}\cos\ang{20}\cos\ang{30}
= 0.2542 + 0.5446 = 0.7987$، إذن $h = \ang{53.0}$.
$\cos\mathrm{A} = (\sin\ang{20} - \sin\ang{48}\sin\ang{53.0})/(\cos\ang{48}
\cos\ang{53.0}) = (0.3420 - 0.5935)/0.4027 = -0.6246$، و$\arccos = \ang{128.7}$؛
فالنجم غرب خط الزوال ($H > 0$)، إذن $\mathrm{A} = \ang{360} -
\ang{128.7} = \ang{231.3}$: أي نحو الجنوب الغربي، على ارتفاع \ang{53} فوق الأفق.
\end{solution}

\begin{solution}{exo:b1:spherical-astronomy:7}
$\cos H_0 = -\tan\ang{60}\tan\ang{23.44} = -1.7321\times0.4335 = -0.7510$،
$H_0 = \ang{138.67}$، فيكون طول النهار $2\times138.67/15 = \qty{18.49}{h}$، أي نحو
\qty{18}{h}~\qty{29}{min}، كما يُظهر منحنى \ang{60}~ش في انقلاب يونيو.
\end{solution}

\begin{solution}{exo:b1:spherical-astronomy:8}
يتقدّم العبور بمقدار \qty{3.93}{min} كل يوم، فيتقدّم بمقدار $30\times3.93 =
\qty{118}{min}$ في 30 يومًا: فيعبر النجم نحو الساعة 20:02. ولكي يعبر عند الظهر يجب أن يتقدّم بمقدار \qty{10}{h} = \qty{600}{min}، وهذا يستغرق
$600/3.93 \approx 153$ يومًا.
\end{solution}

\begin{solution}{exo:b1:spherical-astronomy:9}
(أ) عند أحد الاعتدالين يكون \omterm{def:b1:spherical-astronomy:horizontal}{الارتفاع} الظهري \ang{45}: $I = 1361\times\sin\ang{45}
= \qty{962}{W.m^{-2}}$. (ب) وفي انقلاب يونيو $h = \ang{68.44}$، و
$(a/r)^2 = (149.598/152.100)^2 = 0.9674$:
$I = 1361\times0.9674\times\sin\ang{68.44} = \qty{1224}{W.m^{-2}}$. فالشمس الأعلى تعوّض أكثر مما تنقصه الزيادة في المسافة.
\end{solution}

\begin{solution}{exo:b1:spherical-astronomy:10}
$\sin\delta_\odot = \sin\ang{23.44}\sin\ang{60} = 0.3445$،
$\delta_\odot = \ang{20.15}$. $\tan\alpha_\odot = \cos\ang{23.44}\tan\ang{60}
= 1.5890$، و$\alpha_\odot = \ang{57.82}$ (في الربع نفسه الذي فيه $\lambda_\odot$).
ويتأخر \omterm{def:b1:spherical-astronomy:equatorial}{المطلع المستقيم} بمقدار \ang{2.18}، أي $2.18\times4 \approx
\qty{8.7}{min}$ من الزمن. فالشمس التي تسير بانتظام على \omterm{def:b1:spherical-astronomy:ecliptic}{دائرة البروج} تتحرك مع ذلك بتفاوت في \omterm{def:b1:spherical-astronomy:equatorial}{المطلع المستقيم}، وهو الكمية التي تحدد \omterm{def:b1:spherical-astronomy:equatorial}{الزاوية الساعية} ومن ثم المزولة: فهي متأخرة عن الشمس المتوسطة بين الاعتدال والانقلاب التالي، ومتقدمة عليها بين الانقلاب والاعتدال التالي. وهذا هو حدّ \omterm{def:b1:spherical-astronomy:equatorial}{الميل} في \omterm{def:b1:spherical-astronomy:solar-time}{معادلة الزمن}.
\end{solution}

\begin{solution}{exo:b1:spherical-astronomy:11}
عند الطلوع يكون $h = 0$، فتعطي \cref{prop:b1:spherical-astronomy:transform}
$\sin\delta = \cos\varphi\cos\mathrm{A}_0$؛ وفي انقلاب يونيو
$\delta = \varepsilon$، إذن $\cos\mathrm{A}_0 = \sin\varepsilon/\cos\varphi$.
وعند $\mathrm{A}_0 = \ang{45}$: $\cos\varphi = \sin\ang{23.44}/\cos\ang{45} =
0.5626$، و$\varphi = \ang{55.8}$~ش. وفوق $\ang{90} - \varepsilon =
\ang{66.56}$ يكون $\sin\varepsilon/\cos\varphi > 1$: فلا يوجد سمت للطلوع، لأن الشمس لا تغيب في ذلك اليوم (شمس منتصف الليل).
\end{solution}

\begin{solution}{exo:b1:spherical-astronomy:12}
(أ) الهلالة التي زاويتها $\theta$ هي الكسر $\theta/2\pi$ من الكرة:
$\frac{\theta}{2\pi}\times4\pi R^2 = 2\theta R^2$. (ب) تغطي الأزواج الثلاثة من الهلالات مساحةً مجموعها $2\times(2A + 2B + 2C)R^2$؛ فتغطي الكرة مرة واحدة، إلا المثلث ومقابله فيغطّيان ثلاث مرات، أي بضعف ما يلزم: $4(A + B + C)R^2 = 4\pi R^2 + 4S$، حيث $S$ مساحة المثلث. إذن $S = R^2(A + B + C - \pi)$. (ج) ثلاث زوايا قائمة: $S = R^2(3\pi/2 - \pi) = \pi R^2/2$، أي ثُمن الكرة (ثماني كروي). ومثلث الشكل: $89.9 + 64.9 + 67.4 - 180 = \ang{42.2}$، أي \num{0.736} راديان: $S = 0.736\,R^2$، أي \qty{5.9}{\percent} من الكرة.
\end{solution}

\begin{solution}{pb:b1:spherical-astronomy:1}
\textbf{1.} متممات العرض: $\ang{90} - (\ang{-24.63}) = \ang{114.63}$ و
$\ang{90} - \ang{51.48} = \ang{38.52}$؛ وفرق \omterm{def:b1:spherical-astronomy:geographic}{الطول الجغرافي} \ang{70.40}.

\textbf{2.} $\cos\Delta = \cos\ang{114.63}\cos\ang{38.52} +
\sin\ang{114.63}\sin\ang{38.52}\cos\ang{70.40} = -0.3260 + 0.1899 =
-0.1362$.

\textbf{3.} $\Delta = \ang{97.83}$؛ والمسافة $6371\times1.7074 =
\qty{10878}{km}$، أي نحو \qty{10900}{km}.

\textbf{4.}
\[
  \cos\mathrm{Az} = \frac{\cos\ang{38.52} - \cos\ang{114.63}\cos\ang{97.83}}
  {\sin\ang{114.63}\sin\ang{97.83}} = \frac{0.7823 - 0.0568}{0.9006} = 0.8056 ,
\]
$\mathrm{Az} = \ang{36.3}$، نحو الشمال الشرقي.

\textbf{5.} $\sin\mathrm{Az} = \sin\ang{70.40}\sin\ang{38.52}/\sin\ang{97.83}
= 0.5922$، فيعطي \ang{36.3} أو \ang{143.7}؛ والجيب الموجب في السؤال 4 يختار \ang{36.3}.

\textbf{6.} لا. فالاتجاه الثابت يقطع كل خط زوال بالزاوية نفسها، ويلتف حزونيًا نحو القطب؛ أما اتجاه \omterm{def:b1:spherical-astronomy:spherical-triangle}{الدائرة العظمى} فيتغيّر باستمرار (إنه \ang{36.3} عند البداية وحدها). فالطائرة تقطع مسافة أطول وتصل إلى مكان آخر ما لم تصحّح اتجاهها.

\textbf{7.} \omterm{def:b1:spherical-astronomy:horizontal}{ارتفاع القطب السماوي} الجنوبي يساوي مقدار \omterm{def:b1:spherical-astronomy:geographic}{العرض الجغرافي}: \ang{24.63}.

\textbf{8.} $h = \ang{90} - |\ang{-24.63} - (\ang{-16.72})| = \ang{82.09}$؛
وبما أن $\delta > \varphi$، يعبر سيريوس شمال \omterm{def:b1:spherical-astronomy:celestial-sphere}{سمت الرأس} بمقدار \ang{7.91}.

\textbf{9.} لا يغيب النجم إن كان $\delta < -(\ang{90} - \ang{24.63}) =
\ang{-65.37}$. وكانوبس، عند \ang{-52.70}، ليس دائرًا حول القطب: إنه يغيب.

\textbf{10.} $\ang{90} - |\ang{-24.63} + \ang{52.70}| = \ang{61.93}$، جنوب
\omterm{def:b1:spherical-astronomy:celestial-sphere}{سمت الرأس}.

\textbf{11.} $\cos H_0 = -\tan(\ang{-24.63})\tan(\ang{-16.72}) =
-(0.4584)(0.3004) = -0.1377$؛ $H_0 = \ang{97.92}$. يبقى سيريوس فوق الأفق مدة
$2\times97.92/15 = \qty{13.06}{h}$ من \omterm{def:b1:spherical-astronomy:sidereal-time}{الزمن النجمي}، أي نحو \qty{13}{h}.

\textbf{12.} $\cos\mathrm{A}_0 = \sin(\ang{-16.72})/\cos\ang{24.63} =
-0.3165$: يطلع سيريوس على \omterm{def:b1:spherical-astronomy:horizontal}{السمت} \ang{108.5}، جنوب الشرق.

\textbf{13.} عندما يكون $H = 0$ يكون $\mathrm{LST} = \alpha = \qty{6}{h}~\qty{45}{min}~\qty{9}{s}$.

\textbf{14.} $\qty{86400}{s}\times365.2422/366.2422 = \qty{86164.09}{s} =
\qty{23}{h}~\qty{56}{min}~\qty{4}{s}$؛ يتقدّم العبور بمقدار
\qty{235.9}{s}، أي \qty{3}{min}~\qty{56}{s}، كل ليلة.

\textbf{15.} نحتاج إلى $12 + n\times24/365.2422 = 6.7525 + 24$، إذن
$n = 18.7525/0.065710 = 285.4$ يومًا. و285 يومًا بعد 20 مارس هي 30 ديسمبر: يعبر سيريوس خط الزوال عند منتصف الليل المحلي المتوسط في ليلة 30 إلى 31 ديسمبر 2026.

\textbf{16.} بارانال تقع \ang{70.40} غرب غرينتش: فالزمن المحلي المتوسط متأخر بمقدار $70.40/15 = \qty{4.693}{h}$، أي \qty{4}{h}~\qty{41.6}{min}.

\textbf{17.} \omterm{def:b1:spherical-astronomy:solar-time}{الزمن الشمسي الظاهري} = \omterm{def:b1:spherical-astronomy:solar-time}{الزمن الشمسي المتوسط} + \omterm{def:b1:spherical-astronomy:solar-time}{معادلة الزمن} =
$\qty{0}{h} - \qty{3}{min}$: تُظهر المزولة نحو 23:57.

\textbf{18.} $7\times\qty{3.93}{min} \approx \qty{27.5}{min}$ أبكر.

\textbf{19.} $+\ang{23.44}$ في انقلاب يونيو، و$-\ang{23.44}$ في انقلاب ديسمبر.

\textbf{20.} يونيو: $\ang{90} - |\ang{-24.63} - \ang{23.44}| = \ang{41.93}$.
ديسمبر: $\ang{90} - |\ang{-24.63} + \ang{23.44}| = \ang{88.81}$؛ لأن
$\delta_\odot = \ang{-23.44} > \varphi$، تعبر الشمس شمال \omterm{def:b1:spherical-astronomy:celestial-sphere}{سمت الرأس} بمقدار \ang{1.19}.

\textbf{21.} $I = 1361\times(1/0.9837)^2\times\sin\ang{88.81} = 1361\times
1.0334\times0.9998 = \qty{1406}{W.m^{-2}}$.

\textbf{22.} في يونيو يكون الشتاء الجنوبي. $\cos H_0 = -\tan(\ang{-24.63})
\tan\ang{23.44} = 0.4584\times0.4335 = 0.1988$: تغرب الشمس عند \omterm{def:b1:spherical-astronomy:equatorial}{الزاوية الساعية}
$H_0 = \ang{78.53}$.

\textbf{23.} $\cos\mathrm{A}_0 = \sin\ang{23.44}/\cos\ang{24.63} = 0.4376$:
$\mathrm{A}_0 = \ang{64.0}$، شمال الشرق.

\textbf{24.} طول النهار $2\times78.53/15 = \qty{10.47}{h}$؛ وأطول ليل في السنة في بارانال يدوم $24 - 10.47 = \qty{13.53}{h}$، أي
\textbf{\qty[text-series-to-math=true]{13}{h}~\qty[text-series-to-math=true]{32}{min}}.
\end{solution}
